%% ma: L11-B15-0020
%% lop: 11
%% chuong: 5
%% bai: 15
%% dang: 11.15.3
%% loai: TN4
%% muc: TH
%% dap_an: "C"
%% nguon: "tai-lieu-hoc-tap-mon-toan-lop-11-theo-chuong-trinh-giao-duc-pho-thong-2018.pdf (Lớp Toán Thầy Duy, Phạm Lê Duy, THPT Chu Văn An), Chương 3 Bài 1, Phần 1, Câu 9 (tr.316)"
%% kien_thuc: "Lớp 11 Bài 15 (giới hạn của dãy số)"
%% trang_thai: cho-duyet
%% ly_do: "Dạng mới đề xuất (11.15.x), chờ giáo viên duyệt dạng; tài liệu gốc không có lời giải mẫu, đáp án do kho tự giải và kiểm bằng sympy."
\begin{ex}
Dãy số nào sau đây có giới hạn bằng $0$?
\choice{$u_n=\dfrac{n^2-2}{5n+3n^2}$.}{$u_n=\dfrac{n^2-2n}{5n+3n^2}$.}{\True $u_n=\dfrac{1-2n}{5n+3n^2}$.}{$u_n=\dfrac{1-2n^2}{5n+3n^2}$.}
\loigiai{\sloppy Dãy A và B có giới hạn $\dfrac13$, dãy D có giới hạn $-\dfrac23$; dãy C có bậc tử nhỏ hơn bậc mẫu nên giới hạn bằng $0$.}
\end{ex}
